%% ---------------------------------------------------------------------
\section{扭度 3 与单自旋不对称}\label{sec:power-twist3}

横向单自旋不对称
\begin{equation}\label{eq:power-AN}
A_N\equiv\frac{\dd\sigma(S_{\mathrm T})-\dd\sigma(-S_{\mathrm T})}{\dd\sigma(S_{\mathrm T})+\dd\sigma(-S_{\mathrm T})}
\end{equation}
（$S_{\mathrm T}$ 为一个强子的横向自旋矢量）是幂次修正的一个特别干净的窗口：在领头扭度的共线因子化中它为零。原因有两个。第一，$A_N$ 正比于 $\epsilon^{\mu\nu\rho\sigma}$ 型的“朴素 T 奇”结构（如 $\mathbf S_{\mathrm T}\cdot(\mathbf P_A\times\mathbf P_h)$），要得到实的非零值，振幅之间必须有一个相对相位，即一个吸收部分；第二，横向自旋需要螺旋度翻转，而无质量夸克的手征在微扰论中守恒（与\cref{prop:dis-chiral-suppression}同类的论证）。早在 1978 年就已指出，微扰 QCD 对产生夸克的横向极化的预言在标度极限下为零\cite{Kane:1978nd}。实验上大的 $A_N$（例如 $p^\uparrow p\to\pi X$）因此必须来自领头幂次之外的机制。

当观测到的横动量 $P_{h\mathrm T}$ 远小于 $Q$ 时，相应的机制是 TMD 分布中的 Sivers 函数（第~14~章）；当 $P_{h\mathrm T}\sim Q$ 都是硬标度时，不对称是扭度 3 的效应：强子 $A$ 通过一个夸克—胶子—夸克关联与硬部分相连，所需的相位来自硬部分中某个传播子的极点。描述这一关联的基本对象是 Efremov--Teryaev--Qiu--Sterman（ETQS）函数 $T_F(x_1,x_2)$（\cref{def:tmd-etqs}）\cite{Efremov:1981sh,EfremovTeryaev1985,QiuSterman1991c}，它是带有 $\epsilon^{S_{\mathrm T}\sigma n_An_B}$ 结构的光锥夸克—胶子—夸克矩阵元；其对角值 $T_F(x,x)$ 称为“胶子极点”矩阵元。在两种描述都适用的中间区域 $\Lambda\ll q_{\mathrm T}\ll Q$，二者在单圈层次上给出相同的结果\cite{Ji:2006ub}（\cref{prop:tmd-twist3-matching}）。

\begin{proposition}[SIDIS 型过程中的扭度 3 共线因子化 \tagP]\label{prop:power-twist3-sidis}
在至多一个初态强子的过程中（如 $P_{h\mathrm T}\sim Q$ 处的半单举 DIS），或在两个初态强子而被测末态为色单态、其余末态单举求和的过程中（如 $q_{\mathrm T}\sim Q$ 处的 Drell--Yan）——以下统称“SIDIS 型”——横向单自旋不对称在扭度 3 上具有共线因子化的形式：自旋相关截面 $\dd\Delta\sigma=\frac12[\dd\sigma(S_{\mathrm T})-\dd\sigma(-S_{\mathrm T})]$ 写成硬函数与 $T_F(x,x)$（及其导数）、其他扭度 3 关联函数和未极化部分子分布或碎裂函数的卷积，余项为更高的幂次。具体地：
\begin{enumerate}[label=(\alph*),leftmargin=2em,itemsep=0pt,topsep=2pt]
\item 这一形式在特定过程中已在单圈与次领头阶（NLO）上明确验证：扭度 3 描述与 TMD 描述在重叠区中的一致性在单圈上成立\cite{Ji:2006ub}，Drell--Yan 加权不对称的 NLO 计算见\cite{Vogelsang:2009pj}；
\item $T_F(x,x)$ 的标度演化已经确定：最初几组作者的结果互有出入（\cite{Braun:2009mi}与\cite{Kang:2008ey}），分歧的来源后来被找出，各组结果得到调和\cite{Kang:2012em,Schafer:2012ra}；
\item $T_F(x,x)$ 与 SIDIS 中 Sivers 函数的第一个横动量矩通过算符关系相联系，
\begin{equation}\label{eq:power-TF-sivers}
g\,T_F(x,x)=\mp\int\!\dd^2k_{\mathrm T}\,\frac{\mathbf k_{\mathrm T}^2}{M}\,f^{\perp q}_{1T}(x,k_{\mathrm T})\Big|_{\mathrm{SIDIS}},
\end{equation}
其中符号与归一化依赖于约定（耦合常数的符号、$\epsilon^{0123}$、$T_F$ 的归一化与 Sivers 函数的 Trento 约定）\cite{Boer:2003cm}；本书的约定见\cref{prop:tmd-twist3-matching}。
\end{enumerate}
非微扰地位：\cref{asm:intro-bridge}（MA1、MA2）。
\uses{def:tmd-etqs, prop:tmd-twist3-matching, asm:intro-bridge}
\end{proposition}

\begin{proofsketch}
\begin{enumerate}[label=\arabic*.,leftmargin=1.6em,itemsep=1pt,topsep=2pt]
\item \emph{结构。}把硬部分对强子 $A$ 的部分子横动量展开到一阶，并加上强子另外经一条共线胶子与硬部分相连的图；用 Ward 恒等式组合成规范不变的夸克—胶子—夸克关联（与\cref{prop:power-dis-twist4-x}的共线展开同一方法，但停在一阶）。这一阶的贡献需要一个相位，它来自硬部分中传播子的极点，从而把关联函数“钉”在 $T_F(x,x)$ 或 $T_F(x,x\pm x_h)$ 型的点上（下面的\cref{ex:power-sgp}给出最低阶的明确计算）。
\item \emph{有限阶的验证。}单圈上，TMD 描述（Sivers 函数）在 $q_{\mathrm T}\gg\Lambda$ 处的微扰尾巴与扭度 3 描述在 $q_{\mathrm T}\ll Q$ 处的极限相同\cite{Ji:2006ub}；Drell--Yan 加权单自旋不对称已算到 NLO\cite{Vogelsang:2009pj}。
\item \emph{演化。}$T_F(x,x)$ 的领头阶演化核由几组作者独立推导；在弄清各自对胶子极点处理上的差别之后，结果一致\cite{Braun:2009mi,Kang:2008ey,Kang:2012em,Schafer:2012ra}。
\item \emph{与 Sivers 函数的关系。}\eqref{eq:power-TF-sivers}在算符层面由 Sivers 函数的第一个 $k_{\mathrm T}$ 矩与 Wilson 线结构导出\cite{Boer:2003cm}。在 Kang--Qiu--Vogelsang--Yuan 的约定（$D_\mu=\partial_\mu+\ii gA_\mu$，$\epsilon^{0123}=+1$，Trento 约定）下右边取负号\cite[式~(10)]{Kang:2011hk}；换到本书的 $D_\mu=\partial_\mu-\ii gA_\mu$ 而保持 $T_F$ 的定义式不变时，$gG^{\mu\nu}$ 变号，右边的符号随之反转。该文同时指出这一关系受紫外减除与因子化方案的影响，不是唯一的恒等式，而是一个“零阶”的联系\cite[第~II~节]{Kang:2011hk}。
\end{enumerate}
完整论证见所引文献；不存在全阶论证。
\end{proofsketch}

\paragraph{地位与缺口}
扭度 3 共线因子化的证据都是有限阶的、过程特定的（单圈与 NLO 的明确计算、单圈的匹配、领头阶演化），不存在全阶的扭度 3 论证——既没有次领头幂次的区域分析，也没有积分范数意义下的估计，因此是 \tagP。\eqref{eq:power-TF-sivers}受紫外减除方案的影响\cite[第~II~节]{Kang:2011hk}，其符号又依赖于约定，引用文献中的数值时必须先换算。唯象上，由 SIDIS 拟合的 Sivers 函数的 $k_{\mathrm T}$ 矩与由 $pp$ 中 $A_N$ 拟合的 $T_F(x,x)$ 在大小上大致相当，但符号相反（“符号失配”）\cite{Kang:2011hk}；同时分析 SIDIS、Drell--Yan、$e^+e^-$ 与 $pp$ 数据的全局分析认为各种单自旋不对称有共同的起源，并认为这指向失配的一种解决\cite{Cammarota:2020qcw}。升级途径：次领头幂次（扭度 3）区域分析的全阶论证，以及胶子极点矩阵元的方案无关定义。

\begin{conjecture}[强子—强子碰撞中单举强子单自旋不对称的扭度 3 因子化 \tagC]\label{conj:power-twist3-pp}
强子—强子碰撞 $p^\uparrow p\to hX$ 中大横动量单举强子的单自旋不对称\cite{QiuSterman1991c,Qiu:1998ia}，以及 $p^\uparrow p\to\gamma X$ 中直接光子的单自旋不对称\cite{QiuSterman1992}，在扭度 3 上具有\cref{prop:power-twist3-sidis}那样的共线因子化形式：自旋相关截面等于硬函数与 $T_F(x,x)$、其他扭度 3 关联函数、未极化部分子分布和碎裂函数的卷积，余项为更高的幂次。
非微扰地位：\cref{asm:intro-bridge}（MA1、MA2）。
\uses{prop:power-twist3-sidis, conj:dy-coloured-final, asm:intro-bridge}
\end{conjecture}

\paragraph{地位与缺口}
这里的末态含有大横动量的带色部分子（对直接光子是领头阶子过程 $qg\to\gamma q$、$q\bar q\to\gamma g$ 中反冲的部分子），旁观者与它之间的 Glauber 交换没有被全阶地分析过；即使在领头扭度，带色末态的共线因子化也没有全阶论证（\cref{conj:dy-coloured-final}，\tagC），扭度 3 的陈述继承这一地位。数据方面，见上文的“符号失配”\cite{Kang:2011hk}及其全局分析\cite{Cammarota:2020qcw}；它们是对唯象框架的检验，不是对因子化的证明。升级途径：首先需要领头扭度带色末态因子化的全阶论证。

\begin{example}[最低阶的软胶子极点机制]\label{ex:power-sgp}
考虑\cref{fig:twist3-ssa}中 $A^\uparrow+B\to\gamma+X$ 的 Compton 子过程 $q(x_1P_A)+g(P_B\text{ 侧})\to\gamma+q(p_c)$。强子 $A$ 沿 $+z$ 运动，$P_A=(P_A^+,0^-,\mathbf0_T)$（略去质量）；关联函数 $T_F(x_1,x_2)$ 提供割线两侧的夸克线与一条动量为 $x_gP_A$ 的胶子线，$x_g\equiv x_2-x_1$，后者在振幅一侧接到末态夸克 $c$ 上。夹在硬部分与这条胶子之间的夸克传播子携带 $p_c-x_gP_A$。由于 $p_c^2=0$、$P_A^2=0$，
\begin{equation}\label{eq:power-sgp-den}
(p_c-x_gP_A)^2+\ii0=-2x_g\,p_c\cdot P_A+\ii0=-2\,p_c\cdot P_A\,(x_g-\ii0),\qquad p_c\cdot P_A=p_c^-P_A^+>0,
\end{equation}
最后一步只是把正数 $2p_c\cdot P_A$ 提出，并重新标记无穷小量。用 Sokhotski--Plemelj 公式
\begin{equation}\label{eq:power-sokhotski}
\frac{1}{x_g-\ii0}=\mathrm P\frac{1}{x_g}+\ii\pi\,\delta(x_g),
\end{equation}
传播子分成两部分。主值部分是实的；$\delta$ 函数部分把胶子的动量分数钉在 $x_g=0$（“软胶子极点”：胶子携带零纵向动量，但不是零动量——它的横动量与硬部分耦合），给出因子 $\ii\pi$ 和对角矩阵元 $T_F(x,x)$。正是这个 $\ii$ 与 $T_F$ 中的 $\epsilon^{S_{\mathrm T}\sigma n_An_B}$ 结构相乘，使自旋相关截面成为实的非零量；在树图层次上，除传播子的 $\pm\ii0$ 之外硬部分中没有别的相位来源。

\emph{末态相互作用与初态相互作用。}若胶子不接到出射的夸克 $c$ 上，而是接到来自强子 $B$ 的入射部分子 $b$ 上（Drell--Yan 中的情形），则传播子为 $[(p_b+x_gP_A)^2+\ii0]^{-1}=[2p_b\cdot P_A\,(x_g+\ii0)]^{-1}$。在程函极限下，两种情形给出的耦合因子分别是
\begin{equation}\label{eq:power-sgp-eik}
\bar u(p_c)\,\ii g\gamma^\mu t^a\frac{\ii(\sla p_c-\sla k)}{(p_c-k)^2+\ii0}\;\to\;\bar u(p_c)\,\frac{g\,t^ap_c^\mu}{p_c\cdot k-\ii0},\qquad
\frac{\ii(\sla p_b+\sla k)}{(p_b+k)^2+\ii0}\,\ii g\gamma^\mu t^a u(p_b)\;\to\;-\frac{g\,t^ap_b^\mu}{p_b\cdot k+\ii0}\,u(p_b),
\end{equation}
（$k$ 为流入夸克线的胶子动量，用了 $\bar u(p)\gamma^\mu\sla p=2p^\mu\bar u(p)$、$\sla p\gamma^\mu u(p)=2p^\mu u(p)$；两式与\cref{prop:prelim-wilson-rules}的程函规则一致。）把 \eqref{eq:power-sokhotski}及其共轭代入，可见：相对于各自的主值部分，极点部分在末态相互作用中是 $+\ii\pi\delta$，在初态相互作用中是 $-\ii\pi\delta$。用 Wilson 线的语言，末态相互作用在程函极限下构成未来指向的 Wilson 线，其单胶子项含 $1/(k^+-\ii0)$；初态相互作用构成过去指向的 Wilson 线，含 $1/(k^++\ii0)$；二者之差为 $\frac{1}{k^+-\ii0}-\frac{1}{k^++\ii0}=2\pi\ii\,\delta(k^+)$（\cref{eq:power-sokhotski}，亦即程函规则中“未来减过去”的公式）。因此单自旋不对称在 SIDIS 型（末态相互作用）与 Drell--Yan 型（初态相互作用）过程中符号相反——这正是 Sivers 函数在算符层面的符号反转 $f^{\perp q}_{1T}|_{\mathrm{SIDIS}}=-f^{\perp q}_{1T}|_{\mathrm{DY}}$（\cref{thm:tmd-sivers-sign}，\tagR，$R_{\mathrm{def}}$）在扭度 3 语言中的对应。

\emph{色因子。}在非阿贝尔情形，胶子接到不同的部分子上时色因子也不同；在 $pp\to hX$ 中有多种接法同时出现，且各有不同的色因子与 $\pm\ii0$。这使强子—强子碰撞中的“有效 Sivers 函数”具有比 SIDIS 与 Drell--Yan 之间的简单符号反转更复杂的过程依赖性，也是\cref{conj:power-twist3-pp}只能停留在猜想层次的原因之一。
（\eqref{eq:power-sgp-den}与\eqref{eq:power-sgp-eik}的代数及 \eqref{eq:power-sokhotski}都用符号与数值计算核对过。）
\uses{prop:power-twist3-sidis, thm:tmd-sivers-sign}
\end{example}

%%FIG:fig_twist3_ssa%%
